\section{A specialization to observable states}\label{sec:specialcase}
The distinction between a predictive coordinate, a sampling design, and a reachable direction is explicit in a finite-horizon Markov system with $S_d$ observable states. Let the observation at each stage be the current state. For stages $1,\ldots,H$, use one singleton test for each possible state, all with the empty action string. Write $p^\theta_h(\cdot\mid o,a)$ for a transition distribution. Before the current observation, $q$ is the distribution of that observation given the past. A proper choice of operators for $h<H$ is
\[
 B_h^\theta(o,a)v = v_o\,p^\theta_h(\cdot\mid o,a),
 \qquad B_H^\theta(o,a)v=v_o .
\]
For a signed input, condition on the first state $o$. The remaining policy-weighted absolute suffix mass is $|v_o|$, because after the first state all trajectory probabilities are nonnegative and sum to one. Summing over $o$ yields $\pn{\B^\theta_{H:h}v}=\|v\|_1$. With one action-string group, $\fused v=\|v\|_1$. Thus the class satisfies $\Lambda_B=1$, $U=1$ and $D=S_d$.

The true rank at an interior stage is the dimension of the span of the transition distributions that can be reached there under some preceding state and action. If these distributions are mixtures of $d$ fixed distributions, this rank is at most $d$, even when each distribution assigns positive probability to all $S_d$ observations. For a finite realizable candidate class, the high-probability bounds are $\softO(rAH^4\beta/\eps^2)$ episodes and $\softO(rH)$ calls to each oracle. The history rows still have explicit state coordinates of size $S_d$. In the likelihood budget, a class that independently varies many transition probabilities can have a large $\log|\Theta|$.

For this choice of cores the dual ball is an infinity-norm ball. A coefficient vector $c$ is measured by observing the state and reporting $c_o$, or its normalized version if a unit score is required. A transition coefficient is measured by choosing the prescribed action and then scoring the next state. The general signed-test construction extends these elementary state observations to core tests of several lengths and to policies that branch on intermediate observations. In \sect{sec:example}, both types emit \emph{ready}. The longer stage-two core identifies the probability of type one by imposing mode $a$ and observing the final outcome.
